\section{Separating compatible triples from cancellations}
\label{p13:activity:section}

A count of hosts does not by itself control their probabilities:
compatibility multiplies an independent weight by
\(2^{\rho_{\mathrm{rel}}}\). For incompatible triples,
however, an exact centered identity removes this multiplier.
The remaining cost is a pair coefficient followed by a
dependent target. We estimate that cost before using signs.

\begin{proposition}[The independent cost of one target after a pair]
\label{p13:activity:deficit}
Define
\[
 D_{2,1}=
 \sum_{\substack{e=\{x,y\}\subseteq\mathcal I_N\\d_{xy}>0}}
 |c(e)|
 \sum_{\substack{z\in\mathcal I_N\setminus e\\
            S_z\cap(S_x+S_y)\ne\{0\}}}p_z.
\]
Then \(D_{2,1}\ll_{C,\varepsilon}N^{-3/2+\varepsilon}\).
\end{proposition}

The proof in Section~\ref{p14:comp:activity:deficit} partitions pairs
into finitely many fixed rational channels, weak pairs, and an
exceptional class.  A fixed-source half-count handles the last two
classes.  For a good channel, exact compression turns the pair and its
target into two logarithmic blocks; their total host count is
$N^{1+\varepsilon}$ and their individual weights give an even smaller
$N^{-5/3+\varepsilon}$ contribution.  This bound uses independent
weight only for the new target, while preserving the pair activity.

\begin{theorem}[Absolute activity of incompatible triples]
\label{p13:activity:incompatible}
For every \(\varepsilon>0\),
\[
 \sum_{\substack{T\subseteq\mathcal I_N,\ |T|=3\\q(T)=0}}
           |c(T)|\ll_{C,\varepsilon}N^{-3/2+\varepsilon}.
\]
Both signs of the centered coefficient are included.
\end{theorem}

\begin{proof}
A decomposable triple has a singleton component, hence
centered coefficient zero. For a connected incompatible
triple, expansion of its product gives exactly
\[
 c(T)=-\prod_{i\in T}p_i
       -\sum_{\substack{e\subset T\\|e|=2}}
                  c(e)\prod_{i\in T\setminus e}p_i.
\]
The first terms have total magnitude
\(O(N^{-3/2+\varepsilon'})\) by
Corollary~\ref{p13:half:connected}.
A direct pair has \(c(e)=0\); for any other pair in a
connected triple the remaining space meets its sum
nontrivially. The pair terms are therefore bounded by
three copies of \(D_{2,1}\).
Proposition~\ref{p13:activity:deficit} concludes.
\end{proof}

\begin{proposition}[The third cumulant and the positive activity]
\label{p13:activity:kappa3}
Every compatible triple has \(c(T)\ge0\), and
\begin{equation}
 \kappa_{3,N}=6K_3^++O_{C,\varepsilon}(N^{-3/2+\varepsilon}),
 \qquad
 A_3=K_3^++O_{C,\varepsilon}(N^{-3/2+\varepsilon}).
 \label{p13:activity:kappa-comparison}
\end{equation}
In the second equality the difference \(A_3-K_3^+\)
is nonnegative.
\end{proposition}

\begin{proof}
Order the pair ranks of a compatible triple as
\(a\ge b\ge c\). Their relation spaces show
\(\rho_{\mathrm{rel}}(T)\ge a+b\), so
\[
 \frac{c(T)}{\prod_{x\in T}p_x}
  =2^{\rho_{\mathrm{rel}}(T)}-2^a-2^b-2^c+2\ge0.
\]
If \(b=0\), this is immediate. Otherwise use
\((2^a-1)(2^b-1)\ge2^c-1\).
The same formula includes the primitive case \(a=b=c=0\).
Decomposable triples again have zero coefficient.

In \eqref{p13:exact:kappa3}, the repeated-index terms
are bounded by \(O(p_*A_2+p_*^2\mu_N)
=O(N^{-5/3+\varepsilon'})\).
Apply Theorem~\ref{p13:activity:incompatible} to the
other triples. This proves both formulas.
\end{proof}

For example, three mutually exclusive events have
\(c(T)=2p_xp_yp_z>0\) despite being incompatible.
Thus it is the absolute estimate for incompatible
triples, rather than an assumption on their sign,
that permits the reduction to \(K_3^+\).

\subsection{Pair partitions and the cost of one dependent target}
\label{p14:comp:activity}

The signed triple identities in Section~\ref{p13:activity:section}
use one quantitative input: the independent target deficit following
a pair.  We first prove the finite partition needed for that estimate.

\begin{lemma}[A finite partition of the pair activity]
\label{p13:activity:pair-partition}
For each fixed \(\eta>0\), dependent pairs can be partitioned
into good pairs from a finite fixed set of rational channels,
a weak class, and an exceptional class, with
\[
 A_{\rm w}\ll_{C,\eta,\varepsilon}N^{-1+\eta+\varepsilon},
 \qquad A_{\rm s}\ll_{C,\eta,\varepsilon}N^{-1+\varepsilon}.
\]
A good pair in a channel \(a:b\), \(b\ge3\), has
\[
 |c(e)|\ll_{a,b,C}N^{-(2-1/b)+o(1)}
\]
and an exact compression into a block of fixed logarithmic
width. Moreover \(A_2\ll_{C,\varepsilon}N^{-2/3+\varepsilon}\).
\end{lemma}

\begin{proof}
Choose \(Q\) with \(1/(Q+1)<\eta/2\), and retain the
channels \(3\le b\le Q\), \(b/2<a<b\), \((a,b)=1\).
Lemma~\ref{p1:arith:alignedrank} bounds the rank in all
larger aligned channels by \(L/(Q+1)+o(L)\).
Put these pairs, nonaligned pairs of rank at most a
sufficiently small fixed fraction of \(\eta L\), and
affinely incompatible pairs in the weak class.
Their coefficients have magnitude at most
\(N^{-2+\eta+o(1)}\).
Corollary~\ref{p13:half:pairs} counts their dependent hosts
by \(N^{1+\varepsilon'}\), proving the weak bound.

Nonaligned pairs of macroscopic rank have total activity
\(O(N^{-1+\varepsilon'})\) by Lemma~\ref{p1:arith:dense}.
The local defects contribute only bounded factors under
\eqref{p13:exact:marginals}.
Pairs meeting the irregular set have the same bound:
its marginal mass is \(N^{-1+o(1)}\), by
Proposition~\ref{p1:arith:core}, and fixed-source
propagation \eqref{p13:exact:propagation} bounds their
joint mass. The independent products have the same cost.
At each retained fixed height, normalized defective pairs
cost \(N^{-3/2+1/b+o(1)}\), as in the proof of
Lemma~\ref{p1:arith:fixed-channel-mass}.
Endpoint height-two pairs and non-strict height-one contacts
cost \(N^{-1+o(1)}\), by the same rank and endpoint estimates.
These are the exceptional class.
The remaining pairs have the private coordinates defining
the good fixed-channel compression. Its rank is at most
\(L/b+O(1)\), giving the stated individual weight.

The finite channel sums are \(O(N^{-2/3+\varepsilon'})\),
by Lemma~\ref{p1:arith:fixed-channel-mass}. Together with
the other two classes, for a sufficiently small fixed
\(\eta\), this also proves the assertion about \(A_2\).
\end{proof}

\subsubsection{Proof of Proposition~\ref{p13:activity:deficit}}
\label{p14:comp:activity:deficit}

\begin{proof}
After any fixed pair, Theorem~\ref{p13:half:conditional}
with one target counts at most \(N^{1/2+\varepsilon'}\)
centres. Their total independent mass is at most
\(N^{-1/2+\varepsilon'}\).
The weak and exceptional pair classes thus cost
\(N^{-3/2+\eta+O(\varepsilon')}\).

For a good pair in a fixed channel, multiply its two
windows by the channel coefficients and merge identical
values. An even selection in each original window is
multiplied by a square, so its square class is unchanged.
Merging cancels repeated occurrences and preserves total
even parity. Consequently every nonzero class in
\((S_x+S_y)\cap S_z\) is represented by a nonempty even
selection in the normalized block and in the target.
The normalized block has fixed logarithmic width.
Corollary~\ref{p13:half:pairs} counts these two centres
by \(N^{1+\varepsilon'}\).
The offsets and their inverse normalization have
\(B^{O_{a,b}(1)}\) multiplicity. Each host carries weight
at most \(N^{-(2-1/b)-1+o(1)}\), so its contribution is
\[
 N^{1+\varepsilon'}N^{-(2-1/b)-1+o(1)}
 \ll N^{-5/3+2\varepsilon'}.
\]
Sum over the finite channels, then choose \(\eta\) and
\(\varepsilon'\) within the final margin.
\end{proof}
